\begin{table}[H]

\centering
\caption{Resultados para la \autoref{sec:formulación_de_medios_ejemplo1}}
\label{tab:resultados_formulación_de_medios_ejemplo1}

\begin{threeparttable}
{\small\setstretch{1.35}

\begin{tabular}{M{0.15\linewidth}M{0.15\linewidth}M{0.15\linewidth}M{0.15\linewidth}M{0.15\linewidth}}\hline

    \theader{Compuesto} & \theader{Peso molar} & \theader{g de elemento requeridos} & \theader{g de compuesto requeridos} & \theader{g de compuesto en exceso} \\ \hline
    
    C$_{12}$H$_{22}$O$_{11}$ &    342g/mol & 63.11gC  &            150g           &   187.5g  \\
    (NH$_{4}$)$_{2}$SO$_{4}$ & 132.14g/mol & 18.49gN  &           87.26g          & LIMITANTE\tnote{1} \\
    K$_{3}$PO$_{4}$          & 212.27g/mol &  3.44gP  &  \textbf{23.55g}\tnote{2} &  29.44g   \\
                             &             &  1.28gK  &            2.32g          &  -------- \\
    MgSO$_{3}$               & 104.37g/mol &  1.28gS  &   \textbf{4.17g}\tnote{2} &   5.21g   \\
                             &             &  0.64gMg &            2.78g          &  -------- \\
    CaCl$_{2}$               & 110.98g/mol &  0.38gCa &            1.05g          &   1.31g   \\
    
    \hline
\end{tabular}}

\begin{tablenotes}\tiny % ex: \tntote{1} ... \item[1]
    \item[1] No se calculó el exceso porque es el elemento limitante, y se busca que esté solo en la cantidad justa.; 
    \item[2] Se eligió el valor más alto, pues tal cantidad de compuesto suple ambos elementos.
\end{tablenotes}

\end{threeparttable}
\end{table}